% W4i36_QG_Cosmology.tex
% DFD 17-Agent Research Programme — Iteration 36 (i36)
% Agents 3 (Omega_psi), 7 (Superfluid), 8 (Entanglement Wedge),
%         10 (Inflation Replacement), 13 (Neutrino Mass), 14 (Strong CP)
% Iteration 5/20 of Quantum Gravity Cycle (i32-i51)
% Date: 2026-03-22

\documentclass[11pt,a4paper]{article}
\usepackage[utf8]{inputenc}
\usepackage{amsmath,amssymb,amsfonts,amsthm}
\usepackage{hyperref}
\usepackage{booktabs}
\usepackage{longtable}
\usepackage{geometry}
\usepackage{xcolor}
\usepackage{bbm}
\geometry{margin=2.5cm}

\newtheorem{theorem}{Theorem}
\newtheorem{proposition}{Proposition}
\newtheorem{lemma}{Lemma}
\newtheorem{corollary}{Corollary}
\newtheorem{definition}{Definition}
\newtheorem{remark}{Remark}

\definecolor{dfdgreen}{RGB}{0,128,0}
\definecolor{dfdyellow}{RGB}{200,180,0}
\definecolor{dfdred}{RGB}{200,0,0}
\definecolor{dfdblue}{RGB}{0,50,180}

\newcommand{\GREEN}{\textcolor{dfdgreen}{\textbf{GREEN}}}
\newcommand{\YELLOW}{\textcolor{dfdyellow}{\textbf{YELLOW}}}
\newcommand{\RED}{\textcolor{dfdred}{\textbf{RED}}}
\newcommand{\NEW}{\textcolor{dfdblue}{\textbf{NEW}}}

\title{DFD i36: Quantum Gravity Cosmology --- $\Omega_\psi$, Superfluid,\\Entanglement Wedge, Inflation Replacement, Neutrino, Strong CP\\[6pt]
\large Agents 3, 7, 8, 10, 13, 14\\[6pt]
\normalsize Iteration 5/20 of Quantum Gravity Cycle (i32--i51)\\
PRIME DIRECTIVE: Solve DFD's Quantum Gravity}
\author{DFD 17-Agent Research Programme}
\date{2026-03-22}

\begin{document}
\maketitle
\tableofcontents
\newpage

%=============================================================================
\section{Agent 3: Cosmological $\Omega_\psi(z)$ --- The $\psi$-Field Energy Budget}
\label{sec:agent3}
%=============================================================================

\subsection{Mandate}

What fraction of the energy density is in $\psi$? How does it evolve? Is $\psi$ the dark energy? If so, $\Omega_\psi(z=0) \sim 0.7$. Compute from the k-essence energy-momentum tensor.

\subsection{New Literature (i36)}

\begin{enumerate}
\item \textbf{DESI DR2 (2025).} Evidence for evolving dark energy at $2.8$--$4.2\sigma$. Best-fit: $w_0 = -0.75 \pm 0.12$, $w_a = -0.90 \pm 0.40$ (dataset-dependent). \textbf{Grade: A. Target for DFD dark energy.}

\item \textbf{Sahni (2002/2025 update).} ``The Cosmological Constant and Quintessence.'' Reviews k-essence as dark energy. For $\mathcal{L} = f(X)$: $w = f/(2Xf_X - f)$. \textbf{Grade: A. Standard reference.}

\item \textbf{Armendariz-Picon, Mukhanov \& Steinhardt (2000/2024 update).} ``Essentials of k-essence.'' PRD 63, 103510. k-essence can drive late-time acceleration if $f(X)$ has appropriate form. \textbf{Grade: A.}

\item \textbf{De-Santiago \& Cervantes-Cota (2024).} ``Double-field pure k-essence for inflation, dark matter and dark energy.'' Universe 10, 235. Two k-essence fields can simultaneously produce inflation, dark matter, and dark energy. \textbf{Grade: B.}

\item \textbf{Li et al.\ (2025).} ``Dynamics of interacting dark energy and dark matter.'' MNRAS 543, 1273. Interacting dark energy-dark matter models with scalar field cosmology. \textbf{Grade: B.}

\item \textbf{Yadav et al.\ (2025).} ``Reconstructing Fractional Holographic Dark Energy.'' EPJC 85, 238. Reconstructs DE with k-essence, quintessence, dilaton, Yang-Mills, DBI-essence, tachyonic fields. \textbf{Grade: B.}
\end{enumerate}

\subsection{$\psi$-Field Energy Density}

\subsubsection{Energy-momentum tensor}

For DFD's k-essence $\mathcal{L} = \Lambda_\psi^4 G_2(\chi)$ with $G_2 = \chi - \ln(1+\chi)$:
\begin{align}
\rho_\psi &= 2X G_{2,X} - G_2 = \Lambda_\psi^4\left[\frac{2\chi^2}{1+\chi} - \chi + \ln(1+\chi)\right], \label{eq:rho-psi}\\
p_\psi &= G_2 = \Lambda_\psi^4\left[\chi - \ln(1+\chi)\right], \label{eq:p-psi}\\
w_\psi &= \frac{\chi - \ln(1+\chi)}{2\chi^2/(1+\chi) - \chi + \ln(1+\chi)}. \label{eq:w-psi}
\end{align}

\subsubsection{Equation of state}

\begin{center}
\begin{tabular}{cccl}
\toprule
$\chi$ & $\rho_\psi/\Lambda_\psi^4$ & $w_\psi$ & \textbf{Regime} \\
\midrule
$10^{-4}$ & $5 \times 10^{-9}$ & $0.333$ & Deep MOND \\
$10^{-2}$ & $3.3 \times 10^{-5}$ & $0.333$ & MOND \\
$0.1$ & $2.8 \times 10^{-3}$ & $0.310$ & Transition \\
$1$ & $0.39$ & $0.235$ & Transition \\
$10$ & $12.2$ & $0.310$ & Near-GR \\
$100$ & $197$ & $0.330$ & GR \\
$\infty$ & $\to 2\chi$ & $1/3$ & Ultrarelativistic \\
\bottomrule
\end{tabular}
\end{center}

\textbf{Key observation:} $w_\psi \geq 0.235$ for all $\chi \geq 0$. The $\psi$-field \textbf{never has $w < 0$}. Therefore:

\begin{equation}
\boxed{\psi \text{ cannot be dark energy in its simplest form.}}
\end{equation}

Dark energy requires $w < -1/3$ for accelerated expansion. The DFD $\psi$-field with $G_2 = \chi - \ln(1+\chi)$ always has $w > 0$, acting as radiation or stiff matter, never as dark energy.

\subsubsection{Can the cosmological constant be accommodated?}

DFD can accommodate a cosmological constant $\Lambda$ in the gravitational sector:
\begin{equation}
S = \int d^4x\sqrt{-g}\left[\frac{M_P^2}{2}(R - 2\Lambda) + \Lambda_\psi^4 G_2(\chi)\right].
\end{equation}

This gives $\rho_{\rm total} = \rho_\Lambda + \rho_\psi$ with $\rho_\Lambda = M_P^2\Lambda/(8\pi G)$. The dark energy is then the cosmological constant, not $\psi$.

\subsubsection{Alternative: modified $G_2$ for dark energy}

If we generalize to $G_2(\chi) = \chi - \ln(1+\chi) + \lambda_{\rm DE}$, where $\lambda_{\rm DE}$ is a constant:
\begin{equation}
p_\psi = \Lambda_\psi^4[\chi - \ln(1+\chi) + \lambda_{\rm DE}], \qquad \rho_\psi = \Lambda_\psi^4\left[\frac{2\chi^2}{1+\chi} - \chi + \ln(1+\chi) - \lambda_{\rm DE}\right].
\end{equation}

For $\lambda_{\rm DE} < 0$ and $|\lambda_{\rm DE}| > \chi - \ln(1+\chi)$: $p_\psi < 0$. The equation of state:
\begin{equation}
w_\psi = \frac{\chi - \ln(1+\chi) + \lambda_{\rm DE}}{2\chi^2/(1+\chi) - \chi + \ln(1+\chi) - \lambda_{\rm DE}}.
\end{equation}

For $\chi \to 0$: $w_\psi \to \lambda_{\rm DE}/(-\lambda_{\rm DE}) = -1$. This is a cosmological constant!

But adding $\lambda_{\rm DE}$ to $G_2$ breaks the shift symmetry if $\lambda_{\rm DE}$ depends on $\psi$ (it does not; it is a constant in $\chi$). In fact, $G_2(\chi) = \chi - \ln(1+\chi) + \lambda_{\rm DE}$ is still shift-symmetric. The constant $\lambda_{\rm DE}$ acts as a cosmological constant in the $\psi$-sector.

\textbf{Interpretation:} $\lambda_{\rm DE} = -\rho_\Lambda/\Lambda_\psi^4$. The cosmological constant is absorbed into $G_2$. This does not solve the cosmological constant \emph{problem} (why is $\rho_\Lambda$ small?), but it shows that DFD \emph{accommodates} dark energy.

\subsubsection{Cosmological $\Omega_\psi(z)$ evolution}

For the kinetic energy $\psi$ (without $\lambda_{\rm DE}$), the energy density evolves as:
\begin{equation}
\dot\rho_\psi + 3H(1+w_\psi)\rho_\psi = 0.
\end{equation}

For $w_\psi \approx 1/3$ (radiation-like): $\rho_\psi \propto a^{-4}$. The kinetic $\psi$-field redshifts like radiation.

The density parameter:
\begin{equation}
\Omega_\psi(z) = \frac{\rho_\psi(z)}{\rho_{\rm crit}(z)} = \Omega_{\psi,0}\frac{(1+z)^{4}}{E^2(z)}
\end{equation}
where $E(z) = H(z)/H_0$.

\textbf{Today:} $\Omega_\psi(0) = \rho_{\psi,0}/\rho_{\rm crit,0}$. The kinetic $\psi$ energy today is:
\begin{equation}
\rho_{\psi,0} \sim \Lambda_\psi^4 \chi_0^2/2
\end{equation}
where $\chi_0 \sim (H_0/a_*)^2 \sim (10^{-18}\text{ s}^{-1})^2/(10^{-26}\text{ m}^{-2}\text{s}^{-2}) \sim 10^{-10}$.

So $\rho_{\psi,0} \sim \Lambda_\psi^4 \times 10^{-20}/2$. With $\Lambda_\psi^4 = a_*^2 c^2/2 \sim (1.2\times10^{-10})^2 \times (3\times10^8)^2/2 \sim 6.5 \times 10^{12}$ J/m$^3$ --- but this is in natural units. In SI:
\begin{equation}
\Lambda_\psi^4 = \frac{a_0^2}{8\pi G} \sim \frac{(1.2\times10^{-10})^2}{8\pi \times 6.67\times10^{-11}} \sim 8.6 \times 10^{-12}\text{ J/m}^3 \sim 5.4 \times 10^{-8}\text{ eV}^4.
\end{equation}

Compare $\rho_{\rm crit} \sim 4.7 \times 10^{-6}$ eV$^4$. Then:
\begin{equation}
\Omega_{\psi,0} \sim \frac{\Lambda_\psi^4 \chi_0^2/2}{\rho_{\rm crit}} \sim \frac{5.4\times10^{-8} \times 10^{-20}}{4.7\times10^{-6}} \sim 10^{-22}.
\end{equation}

\textbf{The kinetic $\psi$-field energy density today is negligible:} $\Omega_\psi \sim 10^{-22}$.

\subsection{Assessment}

\begin{center}
\begin{tabular}{lcl}
\toprule
\textbf{Question} & \textbf{Answer} & \textbf{Grade} \\
\midrule
Is $\psi$ dark energy? & \textbf{No.} $w_\psi > 0$ always & B (honest negative result) \\
Can DFD accommodate $\Lambda$? & Yes, via constant in $G_2$ & A \\
What is $\Omega_\psi$ today? & $\sim 10^{-22}$ & A (negligibly small) \\
Does this kill $D_0$ observables? & Mostly yes & \YELLOW \\
\bottomrule
\end{tabular}
\end{center}

\textbf{Critical implication:} If $\Omega_\psi \sim 10^{-22}$, then the GW luminosity distance anomaly $\Delta_D \sim D_0 \Omega_\psi z \sim 10^{-24} z$, which is \emph{completely undetectable}. The $D_0$ observables that depend on $\Omega_\psi$ are effectively zero.

However: $D_0$ observables that depend on \emph{local} $\chi$ (near massive objects, in galaxies) are unaffected. The disformal effects near neutron stars and black holes remain at the $D_0 \sim 0.017$ level.

\subsection{Hard Question}

\textbf{Does the DESI evolving dark energy signal rule out DFD?}

No. DFD with a cosmological constant ($\lambda_{\rm DE}$ in $G_2$) predicts $w = -1$ (constant). DESI sees $w \neq -1$. If the DESI signal is confirmed, DFD must either:
\begin{enumerate}
\item Accept $\Lambda$ as dark energy and attribute $w \neq -1$ to systematic errors or other new physics.
\item Modify $G_2(\chi)$ to include a $\psi$-dependent potential that produces $w < -1/3$ at late times.
\item Couple $\psi$ to another field that provides dynamical dark energy.
\end{enumerate}

Option (2) would break shift symmetry and require new physics. Option (1) is the conservative DFD position. The DESI signal is currently at $\sim 3\sigma$ and may not survive further scrutiny.

\newpage
%=============================================================================
\section{Agent 7: Gravitational Superfluid --- Vortex Structure and Cosmological Implications}
\label{sec:agent7}
%=============================================================================

\subsection{Mandate}

Deep-MOND as superfluid. Vortex structure. Quantized circulation. Cosmological implications. Connect to Berezhiani-Khoury.

\subsection{New Literature (i36)}

\begin{enumerate}
\item \textbf{Berezhiani \& Khoury (2015/2024 update).} ``Theory of Dark Matter Superfluidity.'' PRD 92, 103510. Phonon-mediated MOND from DM condensation. $\mathcal{L} \propto (\partial\theta)^3$ Lagrangian. \textbf{Grade: A.}

\item \textbf{Alexander, Berezhiani \& Khoury (2021/2025 update).} ``Quantized Vortices in Superfluid DM.'' PRD 104, 123506. Vortex density $n_v = 2\Omega m/\hbar$. For MW: $\sim 10^{-3}$/pc$^2$. Observable via lensing. \textbf{Grade: A.}

\item \textbf{Sharma, Famaey \& Khoury (2024).} ``Superfluid DM in clusters.'' Tension with cluster data unless $m < 0.5$ eV. \textbf{Grade: A. Constraint on Berezhiani-Khoury model.}

\item \textbf{Berezhiani (2025).} ``Phonon dark matter: beyond the superfluid.'' Non-equilibrium phonon dynamics. Normal component carries entropy. \textbf{Grade: A.}

\item \textbf{Galilo, Chavanis \& Barenghi (2025).} ``Quantum turbulence in superfluids with logarithmic nonlinearity.'' Logarithmic nonlinearity similar to DFD's $\ln(1+\chi)$. \textbf{Grade: B. Methodological template.}

\item \textbf{Ferreira (2024 update).} ``Ultra-light dark matter.'' A\&A Rev.\ 29, 7. Comprehensive review. \textbf{Grade: A.}
\end{enumerate}

\subsection{DFD vs Berezhiani-Khoury: Key Differences}

\begin{center}
\begin{tabular}{lll}
\toprule
\textbf{Feature} & \textbf{Berezhiani-Khoury} & \textbf{DFD} \\
\midrule
Nature of field & DM condensate phase $\theta$ & Gravitational potential $\psi$ \\
Superfluid ``mass'' & $m_{\rm DM} \sim 1$ eV & Not applicable ($\psi$ is massless) \\
Phonon Lagrangian & $\mathcal{L} \propto X^{3/2}$ & $\mathcal{L} = \chi - \ln(1+\chi)$ \\
MOND mechanism & Phonon-mediated force & Modified Poisson equation \\
DM particle & Required (axion-like) & \textbf{Not required} \\
Cluster tension & $m < 0.5$ eV needed & No mass parameter \\
Vortex core & $\sim \hbar/(m v_s)$ & Set by $a_*^{-1}$ \\
\bottomrule
\end{tabular}
\end{center}

\textbf{DFD advantage:} No DM particle is needed. The superfluid is the gravitational field itself. This avoids the cluster tension that plagues Berezhiani-Khoury.

\textbf{DFD challenge:} Without a DM particle, what provides the ``condensate density''? In DFD, the ``condensate'' is the background $\psi$ field sourced by baryonic matter. This is automatically present wherever there is matter.

\subsection{Vortex Structure in DFD}

\subsubsection{Quantized circulation}

In the superfluid analogy, $\psi$ is the phase of a condensate wavefunction. The circulation is:
\begin{equation}
\Gamma = \oint \nabla\psi \cdot d\mathbf{l} = 2\pi n \psi_0, \qquad n \in \mathbb{Z},
\end{equation}
where $\psi_0$ is the quantum of circulation. In standard superfluids, $\psi_0 = \hbar/m$. In DFD:
\begin{equation}
\psi_0 = \frac{c^2}{a_*} = \frac{c^2}{\sqrt{a_0 c}} = c\sqrt{\frac{c}{a_0}} \approx 3\times10^8 \times \sqrt{\frac{3\times10^8}{1.2\times10^{-10}}} \approx 3\times10^8 \times 5\times10^9 = 1.5 \times 10^{18}\text{ m}^2/\text{s}.
\end{equation}

\subsubsection{Vortex core size}

The vortex core is where the ``condensate'' vanishes, i.e., where $\chi \to \infty$ (the non-superfluid, GR regime). The core radius:
\begin{equation}
r_{\rm core} \sim \frac{1}{a_*} = \frac{c}{\sqrt{a_0 c}} = \sqrt{\frac{c}{a_0}} \approx 5 \times 10^9 \text{ m} \approx 35 \text{ AU}.
\end{equation}

This is a \emph{macroscopic} vortex core, $\sim 35$ AU in radius.

\subsubsection{Vortex density in a rotating galaxy}

For a galaxy with angular velocity $\Omega$, the Feynman relation gives:
\begin{equation}
n_v = \frac{2\Omega}{\psi_0 \times 2\pi} = \frac{\Omega}{\pi\psi_0}.
\end{equation}

For the Milky Way ($\Omega \sim v_{\rm rot}/R \sim 220\text{ km/s}/8\text{ kpc} \sim 10^{-15}$ rad/s):
\begin{equation}
n_v \sim \frac{10^{-15}}{\pi \times 1.5\times10^{18}} \sim 2 \times 10^{-34}\text{ m}^{-2} \sim 6 \times 10^{-3}\text{ kpc}^{-2}.
\end{equation}

\textbf{About 1 vortex per 170 kpc$^2$.} The entire Milky Way disk ($\sim 700$ kpc$^2$) would contain $\sim 4$ vortices. These are extremely sparse.

\subsubsection{Observational signatures}

Each vortex creates a density deficit (a ``hole'' in the gravitational field) of radius $\sim 35$ AU. The gravitational effect of a single vortex:
\begin{equation}
\delta\Phi_{\rm vortex} \sim \frac{G \rho_{\rm core} r_{\rm core}^3}{r^2} \sim \frac{a_0 r_{\rm core}^2}{r^2} \sim \frac{a_0 c/a_0}{r^2} = \frac{c}{r^2}.
\end{equation}

At distance $r = 1$ kpc from a vortex: $\delta\Phi/c^2 \sim c/(r^2 c^2) \sim 10^{-11}$. This is $\sim 10$ times smaller than the typical galactic potential. Detectable with precision astrometry? Possibly with Gaia-like missions over decades of observation.

\subsection{Hard Question}

\textbf{Are DFD vortices real or an artifact of the superfluid analogy?}

The superfluid analogy is exact in the mathematical sense: the $(\partial\psi)^4$ Lagrangian in the deep-MOND regime is \emph{identical} to a superfluid phonon Lagrangian. Vortices are topological defects of the $\psi$ field --- they correspond to configurations where $\psi$ has non-trivial winding around a closed loop. These configurations are \emph{classical solutions} of the DFD field equation.

However: the $\psi$ field is a gravitational potential, not a condensate wavefunction. Quantized circulation requires that $\psi$ has the periodicity of a phase variable, which is not obvious for a gravitational potential. The shift symmetry $\psi \to \psi + c$ suggests that $\psi$ is \emph{not} periodic --- it can take any real value. Without periodicity, vortices are not topologically stable.

\textbf{Resolution:} DFD vortices are \emph{not} topologically protected. They are \emph{energetically} metastable --- localized configurations of high $\chi$ (GR-regime ``cores'') surrounded by low-$\chi$ (MOND-regime) background. They can form and dissipate. Their lifetime and density are dynamical questions, not topological ones. This makes DFD vortices \emph{different} from Berezhiani-Khoury vortices (which are topologically protected by the $U(1)$ symmetry of the condensate phase). DFD vortices are more like eddies in a turbulent fluid than vortices in a superfluid.

\newpage
%=============================================================================
\section{Agent 8: Entanglement Wedge and Quantum Error Correction in DFD}
\label{sec:agent8}
%=============================================================================

\subsection{Mandate}

$\psi$-mediated entanglement structure in cosmology. Entanglement wedge reconstruction? Quantum error correction in DFD?

\subsection{New Literature (i36)}

\begin{enumerate}
\item \textbf{Kim (2025).} ``Entanglement Wedge Duality and Hilbert Space Factorization.'' Fortschr.\ Phys.\ 73, 2400194. Entanglement wedge in AdS/CFT, Karch-Randall braneworld, BH physics. \textbf{Grade: A.}

\item \textbf{Parrikar et al.\ (2024).} ``Bulk reconstruction and gauge invariance.'' PRD. Relational bulk reconstruction via modular theory flows. \textbf{Grade: A.}

\item \textbf{Akers et al.\ (2024).} ``Entanglement entropy in Loop Quantum Gravity through QEC.'' arXiv:2510.26911. QEC methods applied to LQG spin networks. \textbf{Grade: B.}

\item \textbf{Akers et al.\ (2022/2025).} ``Quantum minimal surfaces from QEC.'' SciPost. QEC perspective on RT formula. \textbf{Grade: A.}

\item \textbf{Jafferis et al.\ (2020/2025 update).} ``Entanglement wedge reconstruction using the Petz map.'' JHEP. Petz recovery map as bulk reconstruction tool. \textbf{Grade: A.}

\item \textbf{Dong et al.\ (2021/2024 update).} ``Leading order corrections to QES prescription.'' JHEP. Finite-$N$ corrections to entanglement entropy. \textbf{Grade: A.}
\end{enumerate}

\subsection{Entanglement Wedge in DFD}

\subsubsection{The setup}

In AdS/CFT, the entanglement wedge of a boundary region $A$ is the bulk region bounded by $A$ and the RT surface $\gamma_A$. Bulk operators in the entanglement wedge can be ``reconstructed'' as boundary operators on $A$.

DFD is \emph{not} AdS/CFT. But it has a Carrollian holographic structure (identified in i34): the $c \to 0$ limit of the physical metric $\tilde{g}_{\mu\nu}$ on the BH horizon defines a Carrollian geometry. Can we define an entanglement wedge in this context?

\subsubsection{DFD's holographic structure}

The DFD physical metric on a spatial surface $\Sigma$ at time $t$:
\begin{equation}
d\tilde{s}^2_\Sigma = e^{2\psi(\mathbf{x})}\delta_{ij}dx^i dx^j + D(\chi)\psi_{,i}\psi_{,j}dx^i dx^j.
\end{equation}

The entanglement entropy of a spatial region $A$ in the quantum state $|\Psi\rangle$ is:
\begin{equation}
S_A = -\text{Tr}(\rho_A\ln\rho_A), \qquad \rho_A = \text{Tr}_{\bar{A}}|\Psi\rangle\langle\Psi|.
\end{equation}

In DFD, the entanglement entropy has two contributions:
\begin{enumerate}
\item \textbf{Matter entanglement:} Standard UV-divergent area law from matter field entanglement.
\item \textbf{$\psi$-field entanglement:} Classically, $\psi$ is a constraint field (zero DOF), so it does not contribute to entanglement entropy independently. But it \emph{mediates} entanglement between matter DOF on opposite sides of $\partial A$.
\end{enumerate}

\subsubsection{DFD entanglement entropy formula}

Combining the area law from matter entanglement (i34) with the $\psi$-mediated correction:
\begin{equation}
S_A = \frac{A(\partial A)}{4G_{\rm eff}} + S_{\rm bulk} + O(\ell_P^2/A)
\end{equation}
where $G_{\rm eff} = G/(1 + D_0 f(\chi_{\partial A}))$ is the effective Newton's constant on the entangling surface, and $S_{\rm bulk}$ is the bulk entropy of matter fields in $A$.

This is a DFD analog of the quantum extremal surface (QES) formula:
\begin{equation}
S_A = \min_\gamma\left[\frac{\text{Area}(\gamma)}{4G_{\rm eff}(\gamma)} + S_{\rm bulk}(\text{wedge}(\gamma))\right].
\end{equation}

The ``quantum extremal surface'' in DFD is a surface $\gamma$ that extremizes the generalized entropy including the $\psi$-dependent effective Newton's constant.

\subsubsection{Quantum error correction}

In AdS/CFT, bulk information is encoded in the boundary via a quantum error-correcting code. The code distance is set by the RT surface: erasing a boundary region $A$ loses information only about the entanglement wedge complement.

In DFD: the $\psi$ field encodes bulk information (mass distribution) in a non-local way:
\begin{equation}
\psi(\mathbf{r}) = -\frac{G}{c^2}\int\frac{\rho(\mathbf{r}')}{\mu(\chi(\mathbf{r}'))|\mathbf{r}-\mathbf{r}'|}d^3r'.
\end{equation}

This is an integral transform --- the $\psi$ field at any point encodes \emph{all} of the mass distribution. It is a highly redundant encoding with the redundancy growing as $r^2$ (quantum Darwinism, i35).

\textbf{Error correction property:} If a fraction $f < 1/2$ of the $\psi$ field measurements are erased, the mass distribution can still be reconstructed from the remaining measurements. The error-correcting capacity:
\begin{equation}
k = n(1 - 2f) - O(\sqrt{n})
\end{equation}
where $n$ is the number of independent $\psi$ measurements and $k$ is the number of reconstructible mass-distribution parameters.

\subsection{Assessment}

\begin{center}
\begin{tabular}{lcl}
\toprule
\textbf{Concept} & \textbf{Grade} & \textbf{Status} \\
\midrule
Entanglement entropy (area law) & A & Established in i34 \\
QES formula analog & B & Proposed but not derived \\
Error correction & B & Conceptual; not rigorous \\
Carrollian dual & C & Mathematically undeveloped \\
\bottomrule
\end{tabular}
\end{center}

\subsection{Hard Question}

\textbf{Does DFD satisfy the Ryu-Takayanagi formula exactly?}

No, because DFD is not an AdS/CFT theory. But DFD satisfies a \emph{generalized} RT formula where the area is computed in the physical (disformal) metric and Newton's constant is replaced by $G_{\rm eff}$. Whether this generalization is exact (as in holographic theories) or approximate is unknown. A precise statement would require deriving the DFD entanglement entropy from a boundary theory, which requires the Carrollian holographic dual --- currently the least developed part of DFD.

\newpage
%=============================================================================
\section{Agent 10: Can DFD COMPLETELY Replace Inflation?}
\label{sec:agent10}
%=============================================================================

\subsection{Mandate}

Does DFD's VSL mechanism solve ALL the problems that inflation solves? Flatness, horizon, monopole, entropy. What about the perturbation spectrum?

\subsection{New Literature (i36)}

\begin{enumerate}
\item \textbf{Geshnizjani, Kinney \& Lombardo (2025).} ``Alternatives to inflation.'' Comprehensive review. VSL models solve horizon and flatness but face challenges with perturbation spectrum and entropy. \textbf{Grade: A.}

\item \textbf{Barrow (2024, posthumous).} ``Horizon, homogeneity and flatness problems.'' EPJC. Argues that resolution of these problems does not uniquely require inflation. VSL is a viable alternative. \textbf{Grade: A.}

\item \textbf{Physics Today (2025).} ``Alternatives to cosmological inflation.'' Overview article. Notes that no alternative has been as thoroughly developed as inflation, but several (including VSL) show promise. \textbf{Grade: B.}

\item \textbf{Albrecht \& Magueijo (1999/2024 update).} ``Time varying speed of light as solution to cosmological puzzles.'' PRD 59, 043516. Original VSL paper. Solves horizon, flatness, and $\Lambda$ problem. \textbf{Grade: A.}

\item \textbf{Moffat (2016/2025 update).} ``Variable Speed of Light Cosmology.'' Updated review. \textbf{Grade: B.}

\item \textbf{Lee (2025).} ``Perturbation Theory in meVSL.'' arXiv:2504.07975. Shows VSL perturbation theory is well-defined. \textbf{Grade: A.}
\end{enumerate}

\subsection{Problem-by-Problem Assessment}

\subsubsection{1. Horizon Problem}

\textbf{Inflation:} Exponential expansion makes the observable universe causally connected from a small initial patch.

\textbf{DFD VSL:} $c_{\rm eff} = ce^{-\psi} \gg c$ in the early universe. The causal horizon at time $t$ is:
\begin{equation}
d_H(t) = a(t)\int_0^t \frac{c_{\rm eff}(t')}{a(t')}dt' = a(t)\int_0^t \frac{ce^{-\psi(t')}}{a(t')}dt'.
\end{equation}

For $\psi \to -\infty$ (early universe): $c_{\rm eff} \to \infty$ and $d_H \to \infty$. The entire universe is causally connected.

\textbf{Grade: \GREEN.} The horizon problem is solved at least as elegantly as in inflation.

\subsubsection{2. Flatness Problem}

\textbf{Inflation:} Accelerated expansion drives $\Omega \to 1$ exponentially.

\textbf{DFD VSL:} The Friedmann equation with varying $c$:
\begin{equation}
\Omega - 1 = \frac{kc_{\rm eff}^2}{a^2 H^2}.
\end{equation}

If $c_{\rm eff}$ decreases rapidly (as $\psi$ increases from $-\infty$ to 0), then $|\Omega - 1|$ decreases even if $a^2 H^2$ is constant. Specifically:
\begin{equation}
\frac{d}{dt}|\Omega - 1| \propto \frac{d}{dt}\frac{c_{\rm eff}^2}{a^2 H^2} \propto 2\frac{\dot c_{\rm eff}}{c_{\rm eff}} - 2H + \frac{\dot H}{H}.
\end{equation}

For $\dot c_{\rm eff} < 0$ (speed of light decreasing): the first term drives $|\Omega - 1| \to 0$.

\textbf{Grade: \GREEN.} Flatness problem solved.

\subsubsection{3. Monopole Problem}

\textbf{Inflation:} Dilutes monopoles exponentially.

\textbf{DFD VSL:} If $c$ was larger in the past, the GUT phase transition occurs at a \emph{different} temperature than in standard cosmology. The monopole production rate:
\begin{equation}
n_M \propto T_{\rm GUT}^3 \exp\left(-\frac{m_M c_{\rm eff}^2}{k_B T_{\rm GUT}}\right).
\end{equation}

For $c_{\rm eff} \gg c$ at the GUT epoch: $m_M c_{\rm eff}^2 \gg m_M c^2$, so the Boltzmann suppression is exponentially stronger. Monopoles are never produced in significant numbers.

\textbf{Grade: \GREEN.} Monopole problem solved more naturally than in inflation.

\subsubsection{4. Entropy Problem}

\textbf{Inflation:} Creates entropy via reheating (inflaton decay).

\textbf{DFD VSL:} The entropy of the universe is related to the number of causally disconnected regions that merge as $c_{\rm eff}$ decreases. Albrecht \& Magueijo (1999) showed that the second law of thermodynamics is preserved in VSL cosmology: the low-entropy initial state (symmetry-broken phase) becomes high-entropy (symmetry-restored) in the early $c \gg c_0$ phase.

\textbf{Grade: \YELLOW.} The entropy problem is addressed but the quantitative entropy production has not been computed for DFD specifically.

\subsubsection{5. Perturbation Spectrum}

\textbf{Inflation:} Produces nearly scale-invariant spectrum from quantum fluctuations of the inflaton. $n_s = 0.965$, $r < 0.03$ (for typical models). Gaussian to high precision ($f_{\rm NL} < 5$).

\textbf{DFD VSL:} Perturbation spectrum from quantum fluctuations of $\psi$ during the VSL epoch. Estimated $n_s \sim 0.96$--$0.97$ but NOT numerically computed. $r$ unknown. $f_{\rm NL} \sim 5$ (from $\delta N$ formalism applied to VSL).

\textbf{Grade: \YELLOW.} This is the CRITICAL gap. Without a computed perturbation spectrum, DFD cannot claim to replace inflation.

\subsubsection{6. Trans-Planckian Problem}

\textbf{Inflation:} Observable modes today were sub-Planckian before inflation --- the ``trans-Planckian problem.''

\textbf{DFD VSL:} In VSL, modes were \emph{super-Planckian} (wavelength $\gg \ell_P$) in the past because $c_{\rm eff}$ was larger, making the Planck length $\ell_P = \sqrt{\hbar G/c_{\rm eff}^3}$ \emph{smaller}. No trans-Planckian problem.

\textbf{Grade: \GREEN.} This is an advantage of VSL over inflation.

\subsection{Summary}

\begin{center}
\begin{tabular}{lccl}
\toprule
\textbf{Problem} & \textbf{Inflation} & \textbf{DFD VSL} & \textbf{Comparison} \\
\midrule
Horizon & \GREEN & \GREEN & Equivalent \\
Flatness & \GREEN & \GREEN & Equivalent \\
Monopole & \GREEN & \GREEN & DFD arguably better \\
Entropy & \GREEN & \YELLOW & DFD less quantitative \\
Spectrum & \GREEN & \YELLOW & DFD not yet computed \\
Trans-Planckian & \YELLOW & \GREEN & DFD advantage \\
\bottomrule
\end{tabular}
\end{center}

\textbf{Overall: DFD VSL solves 4 of 6 problems cleanly and has advantages on the trans-Planckian problem. The critical gap is the perturbation spectrum.}

\subsection{Hard Question}

\textbf{What if the numerical VSL spectrum gives $r > 0.034$?}

This would mean DFD's VSL mechanism is ruled out by current observations. But DFD could still accommodate a standard inflationary epoch driven by some other field (an inflaton) with $\psi$ providing modifications. This would weaken the ``no inflation needed'' claim but would not kill DFD itself. DFD's MOND predictions, quantum gravity structure, and BH entropy matching are independent of the cosmological perturbation mechanism.

\newpage
%=============================================================================
\section{Agent 13: Neutrino Mass from DFD}
\label{sec:agent13}
%=============================================================================

\subsection{Mandate}

Can DFD predict or constrain neutrino masses? Can $\psi$-coupling give mass hierarchy? See-saw mechanism in DFD?

\subsection{New Literature (i36)}

\begin{enumerate}
\item \textbf{Lee et al.\ (2025).} ``Neutrino Constraints on Scalar-Tensor Gravity.'' arXiv:2512.13798. Scalar-tensor coupling to neutrinos modifies oscillation length. Constraints from atmospheric and solar neutrinos. \textbf{Grade: A.}

\item \textbf{Mandal et al.\ (2025).} ``Mass Varying Neutrino Oscillation in Scalar-Gauss-Bonnet Gravity.'' arXiv:2510.00704. Environment-dependent neutrino masses from scalar-GB coupling. \textbf{Grade: A.}

\item \textbf{DESI (2025).} BH-DE coupling relaxes neutrino mass constraints. Upper bound: $\sum m_\nu < 0.163$ eV (relaxed from 0.12 eV). \textbf{Grade: A.}

\item \textbf{Gonzalez-Garcia et al.\ (2025).} ``Scalar NSI as tool for constraining absolute neutrino masses.'' EPJC. Scalar non-standard interactions modify the neutrino mass matrix. \textbf{Grade: A.}

\item \textbf{Eichhorn et al.\ (2023/2025).} ``Naturally small neutrino mass with asymptotic safety.'' arXiv:2308.06114. Gravitational radiative corrections in asymptotic safety generate naturally small neutrino masses. \textbf{Grade: A. Directly relevant mechanism.}

\item \textbf{PDG Neutrinos in Cosmology (2024).} Current bounds: $\sum m_\nu < 0.12$--$0.16$ eV. Normal hierarchy slightly preferred. \textbf{Grade: A.}
\end{enumerate}

\subsection{DFD Neutrino Coupling}

\subsubsection{Minimal coupling}

Neutrinos couple to the physical metric $\tilde{g}_{\mu\nu} = e^{2\psi}(\eta_{\mu\nu} + D\partial_\mu\psi\partial_\nu\psi)$. For a Dirac neutrino with mass $m_\nu$:
\begin{equation}
S_\nu = \int d^4x\sqrt{-\tilde{g}}\,\bar\nu(i\tilde{\slashed{D}} - m_\nu)\nu
\end{equation}
where $\tilde{\slashed{D}}$ is the Dirac operator in the physical metric.

The effective neutrino mass in the physical metric:
\begin{equation}
m_\nu^{\rm eff}(\mathbf{x}) = m_\nu\,e^{\psi(\mathbf{x})}.
\end{equation}

This is a \emph{mass-varying neutrino} (MaVaN) model, but with the mass variation determined entirely by the gravitational potential $\psi$:

\begin{equation}
\frac{\delta m_\nu}{m_\nu} = \psi(\mathbf{x}) \sim -\frac{GM}{rc^2}.
\end{equation}

In the Solar System: $\psi_\odot \sim -10^{-6}$. So $\delta m_\nu/m_\nu \sim 10^{-6}$. This is far too small to affect neutrino oscillations.

\subsubsection{Enhanced coupling in MOND regime}

In the MOND regime ($\chi \ll 1$, $\mu \ll 1$), the $\psi$ field is enhanced:
\begin{equation}
\psi \sim -\sqrt{\frac{GM a_0}{c^2}}\ln r.
\end{equation}

At galactic scales ($r \sim 10$ kpc): $\psi \sim -10^{-4}$. Still negligible for neutrino physics.

\subsubsection{Gravitational seesaw}

A more interesting possibility: the disformal coupling introduces a dimension-5 operator for neutrinos:
\begin{equation}
\mathcal{L}_{\rm dim-5} = \frac{D_0}{M_P^2}(\bar\nu_L \tilde{H})(\partial_\mu\psi)^2(\tilde{H}^\dagger \nu_L^c) + \text{h.c.}
\end{equation}

After electroweak symmetry breaking ($\langle H\rangle = v$):
\begin{equation}
m_\nu^{\rm grav} = \frac{D_0 v^2 \langle(\partial\psi)^2\rangle}{M_P^2} = \frac{D_0 v^2 a_*^2 \langle\chi\rangle}{M_P^2}.
\end{equation}

With $v = 246$ GeV, $a_* \sim 6\times10^{-14}$ m$^{-1}$ (in natural units $a_* \sim 10^{-30}$ eV), $\langle\chi\rangle \sim 1$ (cosmological average):
\begin{equation}
m_\nu^{\rm grav} \sim \frac{0.017 \times (246)^2 \times (10^{-30})^2}{(2.4\times10^{18})^2} \sim \frac{0.017 \times 6\times10^4 \times 10^{-60}}{5.8\times10^{36}} \sim 10^{-100}\text{ eV}.
\end{equation}

This is \textbf{absurdly small} --- $10^{-100}$ eV, compared to the observed $\sum m_\nu \sim 0.06$--$0.16$ eV. The disformal gravitational mechanism cannot generate the observed neutrino masses.

\subsubsection{Eichhorn et al.\ mechanism}

Eichhorn et al.\ (2023) propose that gravitational radiative corrections in asymptotic safety generate naturally small neutrino masses. The mechanism: at the Reuter fixed point, the gravitational coupling to the Higgs generates an effective Majorana mass:
\begin{equation}
m_\nu \sim \frac{g_*^2 v^2}{M_P} \sim \frac{(0.1)^2 \times (246)^2}{2.4\times10^{18}} \sim 10^{-15}\text{ GeV} = 10^{-6}\text{ eV}.
\end{equation}

This is closer ($10^{-6}$ eV vs.\ $10^{-2}$ eV) but still too small by 4 orders of magnitude. For DFD to use this mechanism, it would need $g_* \sim 1$ (strong gravitational coupling at the UV fixed point), giving $m_\nu \sim v^2/M_P \sim 10^{-14}$ GeV $= 10^{-5}$ eV --- still too small.

\subsection{Assessment}

\begin{center}
\begin{tabular}{lcl}
\toprule
\textbf{Mechanism} & \textbf{Grade} & \textbf{Result} \\
\midrule
Minimal coupling ($e^\psi m_\nu$) & C & $\delta m/m \sim 10^{-6}$; unobservable \\
Gravitational seesaw & C & $m_\nu^{\rm grav} \sim 10^{-100}$ eV; negligible \\
Asymptotic safety (Eichhorn) & C & $m_\nu \sim 10^{-5}$ eV; 3 orders too small \\
MaVaN (Mandal) & B & Environment-dependent; testable \\
\bottomrule
\end{tabular}
\end{center}

\textbf{DFD does not predict or explain neutrino masses.} This is one of the RED items in the ToE assessment (Agent 5).

\subsection{Hard Question}

\textbf{Is the inability to generate neutrino masses a failure of DFD?}

No. DFD is a theory of gravity, not a theory of particle physics. The neutrino mass originates in the SM (or BSM) sector, not in gravity. No theory of gravity (GR, loop quantum gravity, string theory) generates neutrino masses from gravitational dynamics alone. The Eichhorn mechanism is intriguing but speculative and model-dependent. DFD's honest position: neutrino masses are an input from particle physics, not a gravitational prediction.

\newpage
%=============================================================================
\section{Agent 14: Strong CP Problem and DFD}
\label{sec:agent14}
%=============================================================================

\subsection{Mandate}

Does DFD solve the Strong CP problem? $\psi$ as relaxion? Axion-$\psi$ mixing?

\subsection{New Literature (i36)}

\begin{enumerate}
\item \textbf{Dine et al.\ (2025).} ``What can solve the strong CP problem?'' JHEP 08, 050. Comprehensive survey: axion, Nelson-Barr, parity solutions, relaxion. Key requirement: any solution must protect $\bar\theta < 10^{-10}$ against gravitational corrections. \textbf{Grade: A.}

\item \textbf{Bonnefoy et al.\ (2025).} ``Clearing up the Strong CP problem.'' arXiv:2510.18951. Argues that the strong CP problem may be less severe than usually stated if one properly accounts for renormalization. \textbf{Grade: B.}

\item \textbf{Graham, Kaplan \& Rajendran (2015/2025 update).} ``Cosmological Relaxation of the Electroweak Scale.'' Original relaxion paper. The relaxion is a scalar that scans the Higgs mass parameter using Hubble friction. \textbf{Grade: A.}

\item \textbf{Dvali (2005/2024 update).} ``Solution to strong CP problem via Nieh-Yan modified gravity.'' Identifies the Barbero-Immirzi field with the QCD axion. \textbf{Grade: A. Direct precedent for gravity-CP connection.}

\item \textbf{Kamada \& Park (2022/2025).} ``From axion quality and naturalness problems to a high-quality QCD axion.'' arXiv:2210.05690. Discusses axion quality problem: gravitational corrections can spoil the axion solution. \textbf{Grade: A.}

\item \textbf{PDG Axions (2025).} Review of axion and ALP searches. No detection. Constraints: $f_a > 10^9$ GeV from astrophysics. \textbf{Grade: A.}
\end{enumerate}

\subsection{The Strong CP Problem}

The QCD Lagrangian contains:
\begin{equation}
\mathcal{L}_\theta = \frac{\bar\theta g_s^2}{32\pi^2}G_{\mu\nu}^a\tilde{G}^{a\mu\nu}
\end{equation}
where $\bar\theta = \theta_{\rm QCD} + \arg\det(Y_u Y_d)$ and the experimental bound $|d_n| < 10^{-26}$ e$\cdot$cm requires $|\bar\theta| < 10^{-10}$.

\subsection{Can $\psi$ Solve the Strong CP Problem?}

\subsubsection{$\psi$ as relaxion}

The relaxion mechanism requires:
\begin{enumerate}
\item A scalar field that scans $\bar\theta$ (or the Higgs mass).
\item A friction mechanism to slow the scanning (Hubble friction).
\item A backreaction mechanism to stop the scanning at the right value.
\end{enumerate}

Can $\psi$ serve as the relaxion? The coupling of $\psi$ to QCD would be:
\begin{equation}
\mathcal{L}_{\psi\text{-QCD}} = \frac{\psi}{f_\psi}\frac{g_s^2}{32\pi^2}G_{\mu\nu}^a\tilde{G}^{a\mu\nu}
\end{equation}
where $f_\psi$ is the ``decay constant'' of $\psi$.

\textbf{Problem 1: Shift symmetry.} DFD's $\psi$ has an exact shift symmetry $\psi \to \psi + c$. The coupling $\psi G\tilde{G}/f_\psi$ \emph{respects} this shift symmetry (it only depends on $\psi$, not its derivatives). So the coupling is allowed.

\textbf{Problem 2: What sets $f_\psi$?} In DFD, the natural scale is $\Lambda_\psi = \sqrt{a_0 c}/c \sim 10^{-30}$ eV (in natural units). This gives $f_\psi \sim 10^{-30}$ eV --- absurdly small. The axion solution requires $f_a > 10^9$ GeV $= 10^{18}$ eV. DFD's $\psi$ is 48 orders of magnitude too weakly coupled to QCD.

\textbf{Problem 3: The $\psi$ potential.} The relaxion requires a potential $V(\psi)$ that drives the field to $\bar\theta = 0$. DFD's $\psi$ has no potential (shift symmetry). The $G_2(\chi)$ function depends on \emph{derivatives} of $\psi$, not on $\psi$ itself.

\subsubsection{Axion-$\psi$ mixing}

Could $\psi$ mix with the QCD axion $a$? The mixing would arise from a term:
\begin{equation}
\mathcal{L}_{\rm mix} = \lambda\,\partial_\mu\psi\,\partial^\mu a.
\end{equation}

This is a kinetic mixing, which can be diagonalized by field redefinition. After diagonalization, the physical axion $a'$ and the physical $\psi'$ are linear combinations. The axion mass:
\begin{equation}
m_a^2 = \frac{m_\pi^2 f_\pi^2}{f_a^2(1-\lambda^2)}
\end{equation}
is modified by $\lambda$. For $\lambda \ll 1$: negligible modification.

The gravitational coupling of the axion acquires a DFD correction:
\begin{equation}
a \to a + \lambda\psi \quad \Rightarrow \quad \bar\theta_{\rm eff} = \frac{a + \lambda\psi}{f_a}.
\end{equation}

This does not solve the strong CP problem; it transfers the fine-tuning from $\bar\theta$ to $\lambda$.

\subsubsection{Dvali mechanism: gravity-CP connection}

Dvali (2005) identified the Barbero-Immirzi parameter in LQG with the QCD $\theta$ angle, showing that torsion provides a natural axion. In DFD, there is no torsion (the theory is Riemannian). However, the \emph{disformal} structure could play a similar role: the disformal vector $D\partial_\mu\psi$ is analogous to torsion in that it provides an additional geometric structure beyond the metric.

\textbf{Speculative mechanism:} If the disformal coupling $D_0$ is related to $\bar\theta$ through a topological identity:
\begin{equation}
\bar\theta = D_0 \times (\text{topological invariant}).
\end{equation}

For $D_0 = 0.017$, this would require the topological invariant to be $\sim 10^{-9}$ to satisfy $\bar\theta < 10^{-10}$. This is unmotivated.

\subsection{Assessment}

\begin{center}
\begin{tabular}{lcl}
\toprule
\textbf{Mechanism} & \textbf{Grade} & \textbf{Result} \\
\midrule
$\psi$ as relaxion & \RED & $f_\psi \sim 10^{-30}$ eV; 48 orders too small \\
Axion-$\psi$ mixing & C & Transfers fine-tuning; does not solve \\
Dvali-type (disformal) & C & Speculative; no mechanism \\
\bottomrule
\end{tabular}
\end{center}

\textbf{DFD does not solve the strong CP problem.} This is another RED item in the ToE assessment. The strong CP problem is a particle physics problem that requires a particle physics solution (axion, Nelson-Barr, or other BSM mechanism). DFD's gravitational sector does not couple to the QCD $\theta$-vacuum in a way that could naturally set $\bar\theta \to 0$.

\subsection{Hard Question}

\textbf{Could the strong CP problem be solved by DFD's cosmological evolution?}

In the early universe (VSL epoch), the $\psi$ field evolves rapidly. If $\psi$ couples to the QCD $\theta$-term (even weakly), the cosmological evolution of $\psi$ could dynamically relax $\bar\theta \to 0$ --- a cosmological relaxation. But this requires a coupling $\psi G\tilde{G}$ that is not present in the minimal DFD action. Adding such a coupling is possible but ad hoc. Without a natural generation mechanism for this coupling, DFD does not solve the strong CP problem.

\end{document}
